%! Absolute Value Equations & Inequalities — Unit 1, Lesson 1.2 — Slides (Beamer)
% Deck follows the lesson's gradual-release flow (stats-1.4 spec, 2026-09-03): targets → warm-up →
% hook → I do (notes 1-2, four frames) → we do (guided practice) → debrief → close & start the
% homework. There is no group activity and no individual-practice launch frame: the homework IS the
% individual practice, started in class.
\documentclass[aspectratio=169, 11pt]{beamer}
\usepackage{algebra2-beamer}   % provides \forestheader{} and \sectionlabel[color]{}
% Plain number line: {xmin}{xmax}, ticks labelled every unit.
\newcommand{\numline}[2]{%
  \draw[->, thick, charcoal] (#1-0.5,0)--(#2+0.5,0);
  \foreach \x in {#1,...,#2}{\draw[charcoal] (\x,0.14)--(\x,-0.14) node[below, font=\tiny]{$\x$};}}
% The parkway: mile markers 0..20, ticks every mile, labels every four, tower at mile 12.
\newcommand{\trail}{%
  \draw[->, thick, charcoal] (-0.7,0)--(21.2,0) node[right, font=\tiny]{mile};
  \foreach \x in {0,...,20}{\draw[charcoal] (\x,0.22)--(\x,-0.22);}
  \foreach \x in {0,4,...,20}{\node[below, font=\tiny, charcoal] at (\x,-0.26){$\x$};}
  \fill[forest] (12,0) circle (3pt);
  \node[forest, font=\tiny, above] at (12,0.30){tower};}

\begin{document}

% TITLE SLIDE (CourseName is not defined in beamer, so the name is literal)
{
\setbeamercolor{background canvas}{bg=forest}
\begin{frame}[plain]
  \vspace{2.8cm}\hspace{0.55cm}%
  \begin{minipage}{0.9\textwidth}
    {\color{goldacc}\bfseries\small LESSON 1.2}\\[0.5cm]
    {\color{white}\bfseries\LARGE Absolute Value Equations \& Inequalities}\\[1.2cm]
    {\color{forestmid}\small Algebra 2~$\cdot$~Unit 1: Linear Functions}
  \end{minipage}
\end{frame}
}

% LEARNING TARGETS
\begin{frame}[t, plain]
  \forestheader{Today's targets}
  \sectionlabel{What you will be able to do}
  \begin{itemize}\setlength{\itemsep}{6pt}
    \item Read $|x-h|$ as the \textbf{distance} between $x$ and $h$ --- and say why an equation
          usually has \emph{two} answers.
    \item \textbf{Isolate} the bars, then split into two cases; handle the odd cases where there
          is one answer, or none.
    \item Solve $|X|<c$ and $|X|>c$, and know which gives \emph{one interval} and which gives
          \emph{two rays} --- then write the answer three ways.
  \end{itemize}
  \vspace{4pt}
  \begin{block}{How today runs}
    Warm-up (5) $\to$ notes and one problem together (34) $\to$ debrief (8) $\to$
    \textbf{start the homework, on your own} (13).
  \end{block}
\end{frame}

% WARM-UP
\begin{frame}[t, plain]
  \forestheader{Warm-Up: three things you already know}
  \sectionlabel{5 minutes --- alone, then check with a neighbour}
  \begin{enumerate}\setlength{\itemsep}{7pt}
    \item \textbf{Distance from zero.} $|{-7}|$, \; $|4|$, \; $|{-3}|+|5|$, \; $-|{-6}|$.
          Then: how many different numbers have an absolute value of $7$?
    \item \textbf{Distance between two numbers.} How far apart are $3$ and $11$? Write it with
          bars. And how far is $x$ from $6$?
    \item \textbf{Shade the number line.} $x \ge -1$, and ``between $-3$ and $2$'' --- filled dot
          for included, open dot for not.
  \end{enumerate}
  \vspace{6pt}
  \begin{block}{Hold on to this}
    Both of your pictures came out as \emph{one} connected piece. Could a shaded picture ever come
    in \emph{two} separate pieces? Keep that question.
  \end{block}
\end{frame}

% HOOK
\begin{frame}[t, plain]
  \forestheader{Hook: who can reach the tower?}
  \sectionlabel[goldacc]{A parkway, a radio tower, and 5 miles}
  \begin{center}
  \begin{tikzpicture}[x=0.44cm, y=0.5cm]
    \trail
    \draw[very thick, redacc] (7,0.75)--(17,0.75);
    \fill[redacc] (7,0) circle (3pt);
    \fill[redacc] (17,0) circle (3pt);
    \node[redacc, font=\scriptsize, above] at (9.5,0.8){$5$};
    \node[redacc, font=\scriptsize, above] at (14.5,0.8){$5$};
  \end{tikzpicture}
  \end{center}
  \begin{itemize}\setlength{\itemsep}{4pt}
    \item The tower is at \textbf{mile 12}. A handheld radio reaches it from \textbf{within 5 miles}.
    \item Two hikers stand \emph{exactly} 5 miles out. Where are they? \; \textbf{Miles 7 and 17.}
  \end{itemize}
  \vspace{4pt}
  \begin{block}{The question that carries today}
    \textbf{Why are there \emph{two} of them, and not one?} Every question today is this picture:
    how far is $x$ from the tower --- and is that close enough?
  \end{block}
\end{frame}

% NOTES 1a — DISTANCE
\begin{frame}[t, plain]
  \forestheader{Notes 1: absolute value is a distance --- so it answers twice}
  \sectionlabel{I do --- one centre, two sides}
  \begin{columns}[T]
    \begin{column}{0.52\textwidth}
      \[ |x-h| = \text{the distance between } x \text{ and } h \]
      A hiker at mile $x$ is $|x-12|$ from the tower. ``Exactly 5 miles out'' is $|x-12|=5$:
      \begin{align*}
        x-12 &= 5 \quad\text{or}\quad x-12 = -5 \\
           x &= 17 \quad\text{or}\quad x = 7
      \end{align*}
    \end{column}
    \begin{column}{0.44\textwidth}
      \centering
      \begin{tikzpicture}[x=0.24cm, y=0.44cm]
        \trail
        \draw[very thick, redacc] (7,0.7)--(17,0.7);
        \fill[redacc] (7,0) circle (3pt);
        \fill[redacc] (17,0) circle (3pt);
      \end{tikzpicture}
    \end{column}
  \end{columns}
  \vspace{4pt}
  \textbf{Verify:} $|17-12|=5$ and $|7-12|=5$. \; \textbf{The two-case rule:} $|X|=c$ splits into
  $X=c$ or $X=-c$ --- a distance is never negative, and two spots sit the same distance out.
\end{frame}

% NOTES 1b — ISOLATE
\begin{frame}[t, plain]
  \forestheader{Notes 1, continued: isolate the bars \emph{before} you split}
  \sectionlabel{Undo everything outside the bars first}
  \begin{columns}[T]
    \begin{column}{0.46\textwidth}
      \begin{align*}
        3|x+1| - 2 &= 10 \\
        3|x+1| &= 12 \\
        |x+1| &= 4 \\
        x+1 &= 4 \;\text{ or }\; x+1 = -4 \\
        x &= 3 \;\text{ or }\; x = -5
      \end{align*}
    \end{column}
    \begin{column}{0.50\textwidth}
      \textbf{$|X| = c$ splits into:}
      \begin{itemize}\setlength{\itemsep}{4pt}
        \item $c>0$: \; $X=c$ or $X=-c$ --- \textbf{two}.
        \item $c=0$: \; $X=0$ --- \textbf{one} (only one spot is zero away).
        \item $c<0$: \; \textbf{none} --- a distance is never negative.
      \end{itemize}
    \end{column}
  \end{columns}
  \vspace{5pt}
  \begin{block}{The warning}
    Splitting at the \emph{first} line gives $3x+1-2=10$ --- a single answer, and a wrong one.
  \end{block}
\end{frame}

% NOTES 2a — AND vs OR, three notations
\begin{frame}[t, plain]
  \forestheader{Notes 2: one piece or two --- and three ways to write it}
  \sectionlabel{Same expression, same 5, one symbol apart}
  \begin{columns}[T]
    \begin{column}{0.48\textwidth}
      \textbf{In range: $|x-12| \le 5$} \\[2pt]
      Closer than 5 --- \emph{between}. \\
      $-5 \le x-12 \le 5$, so $7 \le x \le 17$. \\[3pt]
      \centering
      \begin{tikzpicture}[x=0.20cm, y=0.40cm]
        \trail
        \draw[very thick, greenacc] (7,0)--(17,0);
        \fill[greenacc] (7,0) circle (3pt);
        \fill[greenacc] (17,0) circle (3pt);
      \end{tikzpicture} \\[2pt]
      \textbf{One} piece: $\{x \mid 7\le x\le 17\}$, \; $[7,17]$.
    \end{column}
    \begin{column}{0.48\textwidth}
      \textbf{Out of range: $|x-12| \ge 5$} \\[2pt]
      Farther than 5 --- \emph{outside}. \\
      $x-12 \le -5$ \, or \, $x-12 \ge 5$, so $x \le 7$ or $x \ge 17$. \\[3pt]
      \centering
      \begin{tikzpicture}[x=0.20cm, y=0.40cm]
        \trail
        \draw[very thick, redacc] (0,0)--(7,0);
        \draw[very thick, redacc] (17,0)--(20,0);
        \fill[redacc] (7,0) circle (3pt);
        \fill[redacc] (17,0) circle (3pt);
      \end{tikzpicture} \\[2pt]
      \textbf{Two} pieces: $(-\infty,7] \cup [17,\infty)$.
    \end{column}
  \end{columns}
  \vspace{3pt}
  \begin{block}{The endpoint carries the symbol}
    Bracket $[\;]$ + filled dot: \textbf{in} ($\le,\ge$). Parenthesis $(\;)$ + open dot:
    \textbf{out} ($<,>$). $\infty$ \emph{always} takes a parenthesis --- you cannot reach it.
  \end{block}
\end{frame}

% NOTES 2b — THE CRUX
\begin{frame}[t, plain]
  \forestheader{Notes 2, continued: test one hiker and watch the two disagree}
  \sectionlabel[redacc]{The crux --- the wrong symbol hands you the opposite set}
  \begin{center}
  \begin{tikzpicture}[x=0.44cm, y=0.5cm]
    \trail
    \fill[goldacc] (3,0) circle (4pt);
    \node[goldacc, font=\small\bfseries, above] at (3,0.35){mile 3};
    \draw[very thick, greenacc] (7,-0.9)--(17,-0.9);
    \node[greenacc, font=\tiny, below] at (12,-1.0){$|x-12|\le5$};
    \draw[very thick, redacc] (0,-1.7)--(7,-1.7);
    \draw[very thick, redacc] (17,-1.7)--(20,-1.7);
    \node[redacc, font=\tiny, below] at (12,-1.8){$|x-12|\ge5$};
  \end{tikzpicture}
  \end{center}
  \begin{itemize}\setlength{\itemsep}{5pt}
    \item The hiker at mile 3 is $|3-12| = 9$ miles from the tower.
    \item Is $9 \le 5$? \textbf{No.} \quad Is $9 \ge 5$? \textbf{Yes.} \quad Mile 3 is in
          \emph{exactly one} of the two answers --- out of range.
    \item \textbf{The odd cases, same picture:} $|x-12|<-2$ has \textbf{no solution};
          $|x-12|\ge-2$ is \textbf{every real number}. No algebra --- only what a distance is.
  \end{itemize}
\end{frame}

% GUIDED PRACTICE — WE DO
\begin{frame}[t, plain]
  \forestheader{Guided Practice: we do this one together}
  \sectionlabel[goldacc]{You hold the pen --- one problem, three symbols}
  \begin{enumerate}\setlength{\itemsep}{6pt}
    \item[\textbf{a.}] $2|x-5| - 3 = 7$. Isolate, then split. \; Verify one answer.
    \item[\textbf{b.}] $2|x-5| - 3 \le 7$. Closer than or farther than? One piece or two? Write it
          as a solution, an interval, and a shaded line.
    \item[\textbf{c.}] $2|x-5| - 3 \ge 7$. Same three ways.
  \end{enumerate}
  \vspace{5pt}
  \begin{block}{The questions I will ask}
    ``What is stuck to the bars? Undo that first.'' \quad ``Between, or outside?'' \quad
    ``Does your picture have a gap? Then it needs two intervals and a $\cup$.''
  \end{block}
\end{frame}

% DEBRIEF
\begin{frame}[t, plain]
  \forestheader{Debrief: pens down, papers up}
  \sectionlabel[redacc]{8 minutes --- aloud, nothing to fill in}
  \begin{enumerate}\setlength{\itemsep}{5pt}
    \item \textbf{(b) and (c) side by side.} $0\le x\le10$ is \emph{one} piece --- closer than.
          $x\le0$ or $x\ge10$ is \emph{two} --- farther than. Point at the piece, or the gap.
    \item \textbf{The single chain.} $10\le x\le0$ reads as \emph{empty}. Shade it and see why
          (c) needs a $\cup$: $(-\infty,0]\cup[10,\infty)$.
    \item \textbf{Split too early?} $2x-5-3=7$ gives one answer, and a wrong one. Isolate:
          $|x-5|=5$, so $x=10$ or $x=0$.
    \item \textbf{Why never $[\infty$?} You cannot reach it.
  \end{enumerate}
  \vspace{4pt}
  \begin{block}{Say it without the notes}
    The mile-3 hiker: distance 9, not in range, in exactly one of the two answers.
  \end{block}
\end{frame}

% CLOSE & START HOMEWORK
\begin{frame}[t, plain]
  \forestheader{Close --- and start the homework, on your own}
  \sectionlabel{13 minutes $\cdot$ silent $\cdot$ the Homework page in your packet}
  You walked in able to write a line that crosses zero exactly \emph{once}. You walk out able to
  handle a rule that reaches the same value \emph{twice} --- because it measures a \textbf{distance}
  --- and to say whether an answer is one interval or two rays.
  \vspace{5pt}
  \begin{block}{Homework --- scored, due next class}
    We work item~1(a) together right now; then you keep going, alone. Six items: two equations,
    a contrast pair ($<$ against $>$), an inequality to read \emph{off} a shaded number line, the
    odd cases and their boundary, a thermostat, and one multiple-choice item. Whatever is left is
    due next class.
  \end{block}
  \vspace{4pt}
  \textbf{Next:} Lesson 1.3 --- you \emph{solved} $|ax+b|=c$; next you \emph{graph} $y=|ax+b|$, and
  your two answers become the two places a \textbf{V} crosses a horizontal line.
\end{frame}

\end{document}
